% ECONOMICS_PROBLEM: {"id": "I11-632", "title": "Provider payment and patient welfare", "jel_code": "I11", "source_id": "OP-0443"}
% !TeX program = xelatex
\documentclass[11pt,a4paper]{article}
\usepackage[margin=23mm]{geometry}
\usepackage{fontspec}
\setmainfont{Latin Modern Roman}
\usepackage{amsmath,amssymb,mathtools}
\usepackage{xcolor,enumitem,booktabs,longtable,array,xurl,hyperref}
\definecolor{muted}{HTML}{53636C}
\hypersetup{colorlinks=true,urlcolor=blue,linkcolor=blue}
\setlength{\parindent}{0pt}
\setlength{\parskip}{5pt}
\newcommand{\R}{\mathbb R}
\newcommand{\E}{\mathbb E}
\newcommand{\N}{\mathbb N}
\newcommand{\ind}{\mathbf 1}
\newcommand{\Var}{\operatorname{Var}}
\newcommand{\Cov}{\operatorname{Cov}}
\newcommand{\supp}{\operatorname{supp}}
\DeclareMathOperator*{\argmin}{arg\,min}
\DeclareMathOperator*{\argmax}{arg\,max}
\newcommand{\entrymeta}[3]{{\sffamily\small\color{muted}\textbf{\texttt{#1}}\quad|\quad #2\par #3\par}}
\newcommand{\Problemref}[1]{\texttt{#1}}
\newcommand{\JELref}[2]{\texttt{#2} (JEL \texttt{#1})}
\begin{document}
\section*{Provider payment and patient welfare}
\textbf{Problem ID:} \texttt{I11-632}\quad \textbf{JEL:} \texttt{I11}\par
% BEGIN ECONOMICS STATEMENT
\label{problem:OP-0443}
{\sffamily\small\color{muted}Original subject group: Health economics\par}
\label{definitions:problem:30007725}
\textbf{Audit classification:} Published research agenda. Checked source identity and appropriate agenda/background support; expanded intervention, units, population, definedness and causal restrictions. Exact targets and variants are editorial.
\subsection*{Definitions and precise research target}
\textit{Editorial formalization.} Consider primary-care practices serving adults with diabetes in one national health system. Randomize practices to \(a=1\), a specified performance-payment contract, or \(a=0\), their existing payment arrangement, for three years. Let \(Q_i(a)\) denote patient quality-adjusted life-years, \(C_i(a)\) total treatment resource costs, and \(B_i(a)\) appointment and administrative burden. Define
\[\tau=E[\lambda\{Q_i(1)-Q_i(0)\}-\{C_i(1)-C_i(0)\}-\rho\{B_i(1)-B_i(0)\}],\]
where \(\lambda>0\) and \(\rho>0\) are prespecified monetary valuations, and the expectation covers all baseline eligible patients. Measure untreated conditions and avoidable hospitalizations alongside the contract’s rewarded indicators. Under consistency and randomized practice assignment, estimate intention-to-treat effects, adjusting inference for practice clustering. Address spillovers between practices or define a corresponding exposure model. Determine whether improved measured performance yields net patient benefit, and whether effects differ by baseline disadvantage. Transfers to providers enter fiscal accounting, whereas welfare uses their actual resource opportunity cost. This adapts explicit source concerns about incentive effects on quality and unintended outcomes; it does not presume that every performance-pay contract has the same effect.

The practice is the randomization unit and the fixed baseline patient cohort is the target population. Quality-adjusted life-years are the integral or annual sum of a health index \(q\in[0,1]\) multiplied by survival through the three-year horizon; all outcomes use one declared discount and price convention. \(B_i\) is patient hours, \(\rho\) dollars per hour, and \(C_i\) all incremental opportunity costs. Every rewarded metric, payment schedule, risk adjustment, and patient attribution rule must be fixed before assignment. Under no interference, independence of assignment and potential outcomes, positive assignment probabilities, and consistency identify intention-to-treat means; randomization makes this a design/measurement task rather than an unsolved general identification theorem. Program costs outside individual treatment accounting must be added once. Welfare valuations and nonrewarded health outcomes are essential parts of the chosen objective. All expectations use a stated baseline population and probability law. Identification means every admissible data-generating model with the same observed-data law yields the same displayed target; otherwise report the identified set. Exact equations, horizons and assumptions are editorial, rather than source-given canonical conjectures.
\subsection*{Variants and assumption changes}
\begin{enumerate}
\item \textbf{Editorial multitasking variant.} Compare randomized contracts rewarding diabetes alone versus additional prespecified non-diabetes outcomes. Evaluate unrewarded health and net-health benefit.
\item \textbf{Editorial voluntary-entry variant.} Replace assigned payment with a contract offer. Estimate offer effects; receipt effects require instrument assumptions and concern compliers.
\end{enumerate}
\subsection*{References and source audit}
\textbf{1.} Inertia/selection policy questions verified; exact steering experiment and payment estimands are editorial. Locator: Draft 24 March 2014; Section 7, printed p.40; Sections 8–9. Full text or primary publication page inspected. URL: \url{https://kateho.scholar.princeton.edu/sites/g/files/toruqf4136/files/kateho/files/healthcare_io_3_24_14.pdf}.

\textbf{2.} Explicit agenda on broader outcomes, unintended effects, subgroups, and comparative costs. Locator: Authors conclusions; publisher citation/date 1 May 2022. Full text or primary publication page inspected. URL: \url{https://www.cochrane.org/evidence/CD008451_effect-financial-incentives-quality-health-care-provided-primary-care-physicians}.
\begin{enumerate}\raggedright
\item\hypertarget{definitions:source:30007725:1}{} Martin Gaynor; Kate Ho; Robert J. Town. 2014. \emph{The Industrial Organization of Health Care Markets}. Draft 24 March 2014; Section 7, printed p.40; Sections 8–9.
\url{https://kateho.scholar.princeton.edu/sites/g/files/toruqf4136/files/kateho/files/healthcare_io_3_24_14.pdf}.
DOI: \url{https://doi.org/10.3386/w19800}.
\item\hypertarget{definitions:source:30007725:2}{} Anthony Scott; Peter Sivey; Driss Ait Ouakrim; Laura Willenberg; Lena Naccarella; John Furler; Doris Young. 2022. \emph{The effect of financial incentives on the quality of health care provided by primary care physicians}. Authors conclusions; publisher citation/date 1 May 2022.
\url{https://www.cochrane.org/evidence/CD008451_effect-financial-incentives-quality-health-care-provided-primary-care-physicians}.
DOI: \url{https://doi.org/10.1002/14651858.CD008451.pub2}.
\end{enumerate}

\subsection*{Common conventions: Source mathematical conventions}

These conventions apply to each entry unless explicitly replaced.

\textbf{Models and identification.} A primitive model \(m\) specifies all probability laws, feasible choices, information, equilibrium and measurement rules used by its target. The admissible class \(\mathcal M\) is fixed independently of outcomes. If \(P_O(m)\) is its observable probability law and \(\tau(m)\) the target, its identified set at observable law \(P\) is
\[
 \mathcal I_\tau(P)=\{\tau(m):m\in\mathcal M,\ P_O(m)=P\}.
\]
Point identification means that this set is a singleton. An empty set signals incompatibility of data and assumptions. Coordinate bounds need not describe the joint set. Bounds are sharp only if every retained target is attainable or arbitrarily approachable by compatible models. Finite bounds require explicit restrictions on unobserved quantities; unrestricted confounding, hidden wealth, extrapolation or arbitrary equilibrium selection can yield uninformative sets.

\textbf{Causal interventions.} A potential outcome \(Y_i(p)\) is the outcome under a complete policy regime \(p\), including timing, financing, information and market spillovers. The target population and units are fixed before treatment. With interference, use the complete assignment vector or a justified exposure mapping; individual no-interference notation alone cannot identify an economy-wide intervention. A difference of mean potential outcomes does not require the cross-world joint distribution. Conditioning on post-treatment migration, survival, enrollment or employment changes the estimand and needs separate justification.

\textbf{Statistical statements.} An estimator, confidence set, forecast or test specifies a sampling law, model class and loss/coverage criterion. Uniform \((1-\alpha)\)-coverage means \(\inf_{m\in\mathcal M}\Pr_m\{\tau(m)\in C_n\}\geq1-\alpha\), or an explicitly asymptotic version. The whole parameter space trivially has coverage, so informativeness requires a stated width, risk or power objective. A forecasting improvement is relative to a named benchmark, information set and evaluation population.

\textbf{Equilibrium and optimization.} An equilibrium comprises feasible plans, optimality, beliefs where needed, budget constraints and all market/resource clearing. The primitives or a stated benchmark define each correspondence \(\mathcal E(p;m)\). Nonemptiness is assumed or proved, never inferred just from compact policy sets. A selected-equilibrium target uses a declared selection rule; a robust target quantifies over all permitted equilibria. Use suprema or infima unless attainment is established. An \(\varepsilon\)-optimal policy is feasible and within \(\varepsilon\) of the finite optimum value. Report an empty feasible set separately.

\textbf{Welfare and accounting.} Normative utility, interpersonal normalization, social weights, numeraire, discounting and the use of public receipts are primitives. Behavioral utility need not coincide with the welfare criterion. Household welfare includes ownership income and rents once; adding profits again to compensating variation generally double-counts them. A monetary equivalent is defined by an explicit transfer and price convention, with monotonicity and range conditions. Average effects, mechanism contributions, transfers and welfare losses are different targets. Infinite sums require convergence and dynamic feasible plans require terminal or no-Ponzi conditions.

\textbf{Source versions.} A working-paper date, revision date and journal publication year are distinguished. A DOI for a journal version is not automatically the DOI of an earlier manuscript. Locators refer to the stated version and printed pages where specified. A metadata-only check does not verify a claim about a full-text conclusion. Every entry's source subsection narrows the provenance claim accordingly.
% END ECONOMICS STATEMENT
\end{document}
